%% Abstract B – Inverse focal-point synthesis (pdfLaTeX compatible)
%% Overleaf default compiler (pdfLaTeX) works without any changes.
\documentclass[11pt,a4paper]{article}

\usepackage[a4paper, top=1in, bottom=1in, left=1in, right=1in]{geometry}
\usepackage[utf8]{inputenc}
\usepackage[T1]{fontenc}
\usepackage[english]{babel}
\usepackage{newtxtext,newtxmath}   % Times-style body + math fonts
\usepackage{microtype}
\usepackage{setspace}
\setstretch{1.12}
\usepackage{booktabs}
\usepackage[hidelinks]{hyperref}

\setlength{\parindent}{0pt}
\setlength{\parskip}{0.55em}
\widowpenalty=10000
\clubpenalty=10000

\usepackage{titling}
\setlength{\droptitle}{-1.8em}
\pretitle{\begin{center}\large\bfseries}
\posttitle{\par\vskip0.2em\hrule height 0.4pt\par\end{center}\vskip1em}
\preauthor{}\postauthor{}
\predate{}\postdate{}

\begin{document}

\title{Sample-efficient inverse focal-point synthesis for HIFU phased arrays
       via gauge-aware reformulation}
\date{}
\maketitle

\noindent
\textbf{Background.}
Steering a 256-element HIFU phased array onto a desired focus in heterogeneous
tissue is conventionally posed as inverse phase regression, but the map is
doubly ill-posed: empirically $Q\to\text{phase}$ is one-to-many
(nearest-neighbour phase-vector circular-cosine similarity
\textbf{0.002\,$\pm$\,0.043}, indistinguishable from noise) and the system
carries a continuous gauge symmetry---adding a constant $+\delta$ to every
phase produces zero focal displacement and sub-1\,\% intensity change.
Direct deterministic regression cannot fit this combination.

\noindent
\textbf{Method.}
We \textbf{reformulate} the task as \textbf{3-DOF focus-point regression}: a
compact 3-D CNN (FocusPointNet, 0.89\,M params) predicts $(x,y,z)$ from
$(\texttt{log1p}(Q),\,\text{mask})$; transducer phases are then synthesised by
\textbf{analytical delay-and-sum} beamforming, gauge-fixed to $\varphi_0=0$
and quantised to $5^\circ$ steps.
The gauge symmetry was verified three independent ways---analytic derivation
from the delay-and-sum equation, synthetic 256-element least-squares fit, and
full-wave \textbf{k-Wave simulation} (ITU collaborator)---all reporting
\textbf{0\,mm focal shift and $<$1\,\% intensity change at $+20^\circ$}.

\noindent
\textbf{Results.}
On a 22\,/\,4\,/\,4 split with three random seeds, FocusPointNet achieves
\textbf{lateral RMS 4.63\,$\pm$\,2.21\,mm ($X$) and 2.69\,$\pm$\,1.21\,mm
($Y$)}---well inside the 6\,mm focal spot---combined RMS
\textbf{26.25\,$\pm$\,4.61\,mm}, dominated by the elongated focal zone along
the axial axis.
Backbone ablation (CNN\,/\,ResNet-3D\,/\,multi-scale UNet-encoder) and a
state-of-the-art DSNT heatmap variant (nnLandmark\,/\,H3DE-Net 2025) reach
equivalent lateral accuracy: \textbf{precision is a property of the task
formulation, not of any specific backbone}.
Against five closed-form classical baselines on the same split, FocusPointNet
improves lateral $X$ by \textbf{3.0$\times$} and $Y$ by \textbf{5.3$\times$}
over the strongest (amplitude-weighted centroid: 13.82\,/\,14.23\,mm); simpler
peak-finders collapse to 64--80\,mm RMS because the heat peak is offset from
the intended target by 30--50\,mm under refraction.
A $5^\circ$\,/\,$10^\circ$\,/\,$15^\circ$ quantisation study confirms
$<$0.35\,\% intensity error and sub-0.1\,mm focal shift at $5^\circ$, enabling
a discrete gauge-fixed output space.

\noindent
\textbf{Significance.}
To our knowledge this is the first \textbf{gauge-aware reformulation} of HIFU
phased-array inverse synthesis as a low-dimensional regression task, recovering
single-valued ground truth while preserving the underlying physical symmetry.
The empirical gauge-and-quantisation characterisation is of independent value
for classical phase optimisers, and the sample-efficient
learned-localisation\,+\,analytical-beamforming hybrid achieves sub-focal-spot
lateral accuracy from only 22 training simulations.

\smallskip\noindent
\textit{Keywords:} HIFU, phased array, gauge symmetry, focal-point regression,
delay-and-sum, phase quantisation, sample efficiency.

\end{document}
